% 4.5 作为半直积的点群
\section{作为半直积的点群}

本章前几节描述的方法非常普遍，适用于空间群以外的群的分类与约化。特别地，它清楚地揭示了两个点群——其中一个是另一个的子群——的对称性之间的、可恰当地称为\textit{谱系}（genealogical）的关系。

事实是点群都是可解的；也就是说，它们容许一个合成列
\begin{center}
$\mathbf{G} = \mathbf{G}_{1}, \mathbf{G}_{2}, \dots, \mathbf{G}_{m}\ (= E),$
\end{center}
其中 $\mathbf{G}_{i+1}$ 是 $\mathbf{G}_{i}$ 的极大不变子群，而 $\mathbf{G}_{i}/\mathbf{G}_{i+1}$ 是素数阶循环群。（若一个不变子群不包含于另一个阶更大的真不变子群，则称它是极大的。）注意这个可解性的定义乍看之下可能显得太宽松，但对点群它确实是恰当的。

\begin{figure}[H]
\centering
\includegraphics[width=0.66\textwidth]{fig4_1_mu.jpg}
\caption*{\textbf{图 4.1}\ 点群之间的谱系关系。实线表示子群是不变的：(a) 国际记号；(b) Schönflies 记号。}
\end{figure}

现在很容易对所有点群作出半直积，Altmann (1963\textit{a}, \textit{b}) 已考虑过这一问题。我们将把讨论限制在非循环的真转动群上，这是因为其余点群或者是循环群，或者同构于真转动群，或者是直积。表 4.1 以半直积的形式给出点群及其极大不变子群。

\begin{figure}[H]
\centering
\includegraphics[width=0.88\textwidth]{c4_t41_p1.png}
\caption*{\textbf{表 4.1}\ 作为半直积的点群（原书扫描占位）。}
\end{figure}

表 4.1 的注（译文）：

\noindent（i）表的每一行中，依次给出一个群、它的极大不变子群以及相应的半直积。由大括号相连的两行中，第二行是第一行的 Schönflies 版本，符合我们同时给出两种记号的惯例。

\noindent（ii）群名上的撇号表示非标准定向。为获得完整的理解，读者应当用表 2.2 确定表中出现的所有群与子群的元素名称。

\noindent（iii）对 $222\ (D_{2})$，半直积表达式实际上是直积。

\noindent（iv）对 $432\ (O)$，其极大不变子群不是 Abel 群（与其他所有情形不同）。不过，也存在一个不变子群为 Abel 群的表达式，虽然该子群不是极大的，但分解中的另一因子此时是非循环的。

现在举一个例子，说明如何用不变子群的表示来确定半直积群的表示。考虑的情形是 $422 = 4 \wedge 2'\ (D_{4} = C_{4} \wedge C'_{2})$。一个合适的分解是
\begin{equation*}
\begin{array}{l}
422 = E.4 + C_{2x}.4\\
\left(D_{4} = E.C_{4} + C_{2x}.C_{4}\right).
\end{array}
\tag{4.5.1}
\end{equation*}
由表 2.2 可见 $4\ (C_{4})$ 的表示是四个一维表示 $A$、$B$、${}^{1}E$ 和 ${}^{2}E$，其中 $C^{+}_{4z}$ 分别由 $1$、$-1$、$-i$ 和 $i$ 表示。确立 $4\ (C_{4})$ 为不变子群的共轭关系是
\begin{center}
$C^{-1}_{2x}EC_{2x} = E$，\quad $C^{-1}_{2x}C^{+}_{4z}C_{2x} = C^{-}_{4z}$，\quad $C^{-1}_{2x}C_{2z}C_{2x} = C_{2z}$，\quad 以及 $C^{-1}_{2x}C^{-}_{4z}C_{2x} = C^{+}_{4z}$。
\end{center}
$4\ (C_{4})$ 的表示分成三颗星：单由 $A$ 组成的、单由 $B$ 组成的、以及由 ${}^{1}E$ 与 ${}^{2}E$ 合起来组成的。${}^{1}E$ 与 ${}^{2}E$ 走到一起，源于形如 ${}^{1}\mathbf{E}(C^{-1}_{2x}C^{+}_{4z}C_{2x}) = {}^{1}\mathbf{E}(C^{-}_{4z}) = {}^{2}\mathbf{E}(C^{+}_{4z})$ 的方程。

三种情形下的小陪群是：对星 $A$ 和 $B$，是含 $E$ 与 $C_{2x}$ 的群 $2'\ (C'_{2})$；对其余的星 $E$，是群 $1\ (C_{1})$。由于这是半直积，小陪群所需的表示就是它们的普通表示，这些可以从表 2.2 得到。

设 $\phi$ 是 ${}^{1}E$ 的基函数，于是
\begin{equation*}
C^{+}_{4z}\phi = \phi\,{}^{1}\mathbf{E}(C^{+}_{4z}) = -\mathrm{i}\phi.
\tag{4.5.2}
\end{equation*}
记 $C_{2x}\phi = \psi$，则 $\langle \phi, \psi |$ 是表示 ${}^{1}E \uparrow 422\ (D_{4}) = \Delta$ 的一个基。我们计算 $C^{+}_{4z}$ 与 $C_{2x}$ 的诱导矩阵表示。现在
\begin{equation*}
C^{+}_{4z}\psi = C^{+}_{4z}C_{2x}\phi = C_{2x}C^{-}_{4z}\phi = C_{2x}\phi\,{}^{1}\mathbf{E}(C^{-}_{4z}) = i\psi.
\tag{4.5.3}
\end{equation*}
由式 (4.5.2) 与 (4.5.3) 可见
\begin{equation*}
C^{+}_{4z}\langle \phi, \psi | = \langle -\mathrm{i}\phi, \mathrm{i}\psi | = \langle \phi, \psi |\begin{pmatrix} -\mathrm{i} & 0\\ 0 & \mathrm{i} \end{pmatrix}
\end{equation*}
因此
\begin{equation*}
\boldsymbol{\Delta}(C^{+}_{4z}) = \begin{pmatrix} -\mathrm{i} & 0\\ 0 & \mathrm{i} \end{pmatrix}.
\tag{4.5.4}
\end{equation*}
另外，按定义
\begin{equation*}
C_{2x}\phi = \psi
\tag{4.5.5}
\end{equation*}
于是
\begin{equation*}
C_{2x}\psi = C_{2x}C_{2x}\phi = \phi,
\tag{4.5.6}
\end{equation*}
从而
\begin{equation*}
\Delta(C_{2x}) = \begin{pmatrix} 0 & 1\\ 1 & 0 \end{pmatrix}.
\tag{4.5.7}
\end{equation*}
上述结果汇总于表 4.2；为了与表 2.2 一致，我们作认同 $\Lambda^{A}_{1} = A_{1}$、$\Lambda^{A}_{2} = A_{2}$、$\Lambda^{B}_{1} = B_{1}$、$\Lambda^{B}_{2} = B_{2}$、$\Delta = E$。

\begin{table}[H]
\centering
\caption*{\textbf{表 4.2}\ $422\ (D_{4})$ 的表示与 $4\ (C_{4})$ 的表示的关系（Altmann 1963\textit{a}）}
\begin{tabular}{@{}lllllllllllll@{}}
\toprule
$4\ (C_{4})$ & 星 & 小陪群 & $E$ & $C^{+}_{4z}$ & $C_{2z}$ & $C^{-}_{4z}$ & $422\ (D_{4})$ & $E$ & $C^{+}_{4z}$ & $C_{2z}$ & $C_{2x}$ & $C_{2a}$\\
 & & & & & & & & & $C^{-}_{4z}$ & & $C_{2y}$ & $C_{2b}$\\
\midrule
$A$ & $A$ & $2'\ (C'_{2})$ & 1 & 1 & 1 & 1 & $A_{1}$ & 1 & 1 & 1 & 1 & 1\\
 & & & & & & & $A_{2}$ & 1 & 1 & 1 & $-1$ & $-1$\\
$B$ & $B$ & $2'\ (C'_{2})$ & 1 & $-1$ & 1 & $-1$ & $B_{1}$ & 1 & $-1$ & 1 & 1 & $-1$\\
 & & & & & & & $B_{2}$ & 1 & $-1$ & 1 & $-1$ & 1\\
${}^{1}E$ & $E$ & $1\ (C_{1})$ & $\begin{cases}1\\ 1\end{cases}$ & $\begin{cases}-i\\ i\end{cases}$ & $\begin{cases}-1\\ -1\end{cases}$ & $\begin{cases}i\\ -i\end{cases}$ & $E$ & 2 & 0 & $-2$ & 0 & 0\\
${}^{2}E$ & & & & & & & & & & & &\\
\bottomrule
\end{tabular}
\end{table}

表 4.3 复制了表 4.1 中出现的其余半直积的类似关系的汇总（原书扫描占位）。

\begin{figure}[H]
\centering
\includegraphics[width=0.92\textwidth]{c4_t43_p1.png}
\caption*{\textbf{表 4.3}\ 点群的表示与其不变子群的表示的关系（Altmann 1963\textit{a}；原书扫描占位）。}
\end{figure}
\begin{figure}[H]
\centering
\includegraphics[width=0.92\textwidth]{c4_t43_p2.png}
\caption*{\textbf{表 4.3}（第 2 页）。}
\end{figure}
\begin{figure}[H]
\centering
\includegraphics[width=0.92\textwidth]{c4_t43_p3.png}
\caption*{\textbf{表 4.3}（第 3 页）。}
\end{figure}
\begin{figure}[H]
\centering
\includegraphics[width=0.92\textwidth]{c4_t43_p4.png}
\caption*{\textbf{表 4.3}（第 4 页）。}
\end{figure}

在表 4.3 的例子中，只有分解 $432 = 222 \wedge 32'\ (O = D_{2} \wedge D'_{3})$ 真正有些困难。因为它提供了这样一个例子：小群既不是整个群也不是不变子群。在星 $B$ 中，小群是 $422'\ (D'_{4})$（小群的取向依赖于从 $B_{1}$、$B_{2}$ 还是 $B_{3}$ 出发），而小表示恰是 $422'\ (D'_{4})$ 的两个一维表示；这就是为什么我们得到 $432\ (O)$ 的两个三维表示——$432\ (O)$ 的阶是 $422'\ (D'_{4})$ 阶的三倍。空间群理论中相应的情形是：所考虑的 $\mathbf{k}$ 空间的点具有比平移群更多的对称性，但又不具有空间群的全部对称性。

有趣的是，把四元数群 $\mathbf{Q}$（表 5.1 中的群 $\mathbf{G}^{5}_{8}$）的表示——由 4 阶循环群（由于不把它当作点群，记作 $\mathbf{G}^{1}_{4}$）出发构成——的生成方式，与上面从同一个循环群生成 $422\ (D_{4})$ 的表示的方式加以对比。设 $\mathbf{G}^{1}_{4}$（4 阶循环群）的元素是 $E, P, P^{2}, P^{3}$（考虑 $422\ (D_{4})$ 时 $P$ 对应于 $C^{+}_{4z}$），则四元数群有 8 个元素，分为 5 个类：$E$；$P^{2}$；$P, P^{3}$；$Q, P^{2}Q$；以及 $PQ, P^{3}Q$。也可以直接说生成关系是 $P^{4} = E$、$Q^{4} = E$、$P^{2} = Q^{2}$、$QP = P^{3}Q$。左陪集分解为
\begin{equation*}
\mathbf{Q} = E\mathbf{G}^{1}_{4} + Q\mathbf{G}^{1}_{4}.
\tag{4.5.8}
\end{equation*}
$Q\mathbf{G}^{1}_{4}$ 的所有元素都是 4 阶的，所以群 $\mathbf{Q}$ 不是半直积。它在空间群理论中类似于一个非点式空间群。与之前一样，这里有三颗星 $A$、$B$ 和 $E$，小群分别是 $\mathbf{Q}$、$\mathbf{Q}$ 和 $\mathbf{G}^{1}_{4}$。星 $E$ 同样包含两个表示 ${}^{1}E$ 与 ${}^{2}E$。小陪群是：对星 $A$ 和 $B$，是我们记作 $\mathbf{G}^{1}_{2}$ 的 2 阶循环群；对星 $E$，是仅由恒等元组成的平凡群 $\mathbf{G}^{1}_{1}$。由于此情形不再是半直积，小陪群所需的表示需要进一步考察。对由元素 $e$ 与 $q$ 组成的小陪群 $\mathbf{G}^{1}_{2}$，我们需要（可能是射影的）表示使得
\begin{equation*}
\boldsymbol{\Lambda}^{A}(q)\boldsymbol{\Lambda}^{A}(q) = \mathbf{A}(P^{2})\boldsymbol{\Lambda}^{A}(e),
\tag{4.5.9}
\end{equation*}
以及
\begin{equation*}
\boldsymbol{\Lambda}^{B}(q)\boldsymbol{\Lambda}^{B}(q) = \mathbf{B}(P^{2})\boldsymbol{\Lambda}^{B}(e).
\tag{4.5.10}
\end{equation*}
然而 $\mathbf{A}(P^{2}) = \mathbf{B}(P^{2}) = 1$，所以小表示再一次是普通矢量表示。于是四个一维表示 $\Lambda^{A}_{1}, \Lambda^{A}_{2}, \Lambda^{B}_{1}, \Lambda^{B}_{2}$ 对 $\mathbf{Q}$ 具有与对 $422\ (D_{4})$ 完全相同的结构。对星 $E$，我们从表示 ${}^{1}E$ 构造表示 $\Gamma = {}^{1}E \uparrow \mathbf{Q}$，这是从星 $E$ 得到的 $\mathbf{Q}$ 的唯一表示。记 ${}^{1}E$ 的基函数为 $\phi$，使得
\begin{equation*}
P\phi = -\mathrm{i}\phi,
\tag{4.5.11}
\end{equation*}
并记 $Q\phi = \psi$，则 $\langle \phi, \psi |$ 是 $\Gamma$ 的一个基。现在
\begin{equation*}
P\psi = PQ\phi = QP^{3}\phi = \mathrm{i}Q\phi = \mathrm{i}\psi
\tag{4.5.12}
\end{equation*}
于是由式 (4.5.11) 与 (4.5.12) 可见
\begin{equation*}
\Gamma(P) = \begin{pmatrix} -\mathrm{i} & 0\\ 0 & \mathrm{i} \end{pmatrix}.
\tag{4.5.13}
\end{equation*}
这与 $422\ (D_{4})$ 的情形完全一样（见式 (4.5.4)）。另外 $Q\phi = \psi$（按定义），于是
\begin{equation*}
Q\psi = Q^{2}\phi = P^{2}\phi = -\phi.
\tag{4.5.14}
\end{equation*}
这里终于出现了 $\mathbf{Q}$ 与 $422\ (D_{4})$ 的表示之间的唯一差别：现在
\begin{equation*}
\boldsymbol{\Gamma}(Q) = \begin{pmatrix} 0 & -1\\ 1 & 0 \end{pmatrix}
\tag{4.5.15}
\end{equation*}
这与 $\Delta(C_{2x})$ 的式 (4.5.7) 形成对照。但尽管如此，两个群的特征标表是相同的。这展示了一个显著的事实：两个不同构的群可以有完全相同的特征标表。

前面已经提到，每个可解的有限群 $\mathbf{G}$ 都有素数指数的极大不变子群 $\mathbf{N}$。特别地，每个晶体学点群都含有一个指数为 2 或 3 的极大不变子群。由此可以证明每个空间群也含有指数为 2 或 3 的极大不变子群（Zak 1960）。因此，作为 4.1--4.4 节理论的特殊情形，考察当 $\mathbf{G}/\mathbf{N}$ 是素数 $p$ 阶循环群时 $\mathbf{G}$ 的表示如何从 $\mathbf{N}$ 的表示导出，是合适的。这一理论构成了 Seitz (1936\textit{b}) 原本给出的、说明如何约化空间群的基础；近年 Zak (1960) 采用的某些方法也依赖于此理论。当然，4.4 节的理论对不变子群为 Abel 群的情形（见该节情形 (i)）已覆盖得很好；这对空间群已相当足够，因为其不变子群可取为平移不变子群 $\mathbf{T}$。然而下面包括的理论不仅仅作为数学练习，也不仅仅由于其历史意义，而是因为它进一步揭示了我们处理的这些群的结构。例如，它说明了如何利用像 $432 = 23 \wedge 2''\ (O = T \wedge C''_{2})$（见表 4.1）这样的分解——由于 $23\ (T)$ 不是 Abel 群，它不落在 4.4 节情形 (i) 的范围内。附带地，该理论表明当 $\mathbf{G}/\mathbf{N}$ 为素数阶时，4.4 节的情形 (ii) 并不像乍看之下那样难于处理。

设 $\mathbf{G}$ 是群，$\mathbf{N}$ 是 $\mathbf{G}$ 的素数指数 $p$ 的极大不变子群。此时可以把 $\mathbf{G}$ 分解成左陪集
\begin{equation*}
\mathbf{G} = \mathbf{N} + r\mathbf{N} + r^{2}\mathbf{N} + \dots + r^{p-1}\mathbf{N}
\tag{4.5.16}
\end{equation*}
其中 $r^{p} \in \mathbf{N}$。当 $r^{p}$ 等于恒等元 $E$ 时，$\mathbf{G}$ 是 $\mathbf{N}$ 与 $p$ 阶循环群的半直积。更一般地，设 $r^{p} = \bar{n} \in \mathbf{N}$ 且 $\bar{n}$ 的阶为 $l$。则 $E = \bar{n}^{l} = r^{lp}$。若 $1 \leqslant l \leqslant p$，则由于 $p$ 是素数，可以选 $s = r^{l}$ 代替 $r$，组成另一种陪集分解
\begin{equation*}
\mathbf{G} = \mathbf{N} + s\mathbf{N} + s^{2}\mathbf{N} + \dots + s^{p-1}\mathbf{N}
\tag{4.5.17}
\end{equation*}
此时 $s^{p} = E$。于是取上述范围内的 $l$ 值，就可以重新选取 $r$ 使 $l = 1$。另一方面，若 $l = p$，则不存在这种替代选择。沿这些路线的进一步论证表明，总可以这样选取 $r$，使得 $l$ 取 $p$ 的幂值，即 $l = 1, p, p^{2}, \dots$。因此在分解 (4.5.16) 中可以不失一般性地假定 $r^{(p^{m})} = E$，其中 $m \geqslant 1$。

接下来考察 $\mathbf{G}$ 相对于 $\mathbf{N}$ 的类结构。设 $C_{1}, C_{2}, \dots, C_{n}$ 是 $\mathbf{N}$ 的类。可以证明用 $r$ 作共轭会给出这些类的一个置换。设 $n_{i} \in C_{i}$ 且 $rn_{i}r^{-1} = n_{j} \in C_{j}$。设 $n'_{i}$ 是 $C_{i}$ 的任一其他元素，则存在 $n \in \mathbf{N}$ 使 $n'_{i} = nn_{i}n^{-1}$，且 $rn'_{i}r^{-1} = rnn_{i}n^{-1}r^{-1}$。由于 $\mathbf{N}$ 在 $\mathbf{G}$ 中不变，存在 $n' \in \mathbf{N}$ 使 $rn = n'r$。于是 $rn'_{i}r^{-1} = n'(rn_{i}r^{-1})n'^{-1} = n'n_{j}n'^{-1} \in C_{j}$。这表明 $rC_{i}r^{-1} \subseteq C_{j}$，从而 $C_{i}$ 的元素数 $s_{i}$ 小于或等于 $C_{j}$ 的元素数 $s_{j}$。用 $r^{-1}$ 作共轭的类似论证给出 $r^{-1}C_{j}r \subseteq C_{i}$ 且 $s_{j} \leqslant s_{i}$。因此 $rC_{i}r^{-1} = C_{j}$ 且 $s_{i} = s_{j}$。设 $C_{k}$ 是另一个使 $rC_{k}r^{-1} = C_{j}$ 的类，则 $C_{i} = r^{-1}C_{j}r = C_{k}$。故用 $r$ 作共轭产生 $\mathbf{N}$ 的类上的一个置换 $\pi(r)$。类似地可以证明，用 $r^{m}$ 作共轭产生 $\mathbf{N}$ 的类上的置换 $\pi(r^{m})$（$m = 1, 2, \dots p$）。由于 $r^{p} \in \mathbf{N}$，$\pi(r^{p})$ 是恒等置换 $\varepsilon$。又有 $\{\pi(r)\}^{2}C_{i} = \pi(r)rC_{i}r^{-1} = r^{2}C_{i}r^{-2} = \pi(r^{2})C_{i}$，故 $\{\pi(r)\}^{2} = \pi(r^{2})$，类似地 $\{\pi(r)\}^{m} = \pi(r^{m})$。于是 $\{\pi(r)\}^{p} = \pi(r^{p}) = \varepsilon$。因此 $\pi(r)$ 的阶整除 $p$。由于 $p$ 是素数，$\pi(r)$ 的阶为 1 或 $p$。这意味着把 $\pi(r)$ 分解为不相交循环的乘积时，任何循环的长度是 1 或 $p$。所以 $\mathbf{N}$ 的类有两种：在 $r$ 下自共轭的类（$rC_{i}r^{-1} = C_{i}$），以及以长度为 $p$ 的循环出现、使 $C_{i}, rC_{i}r^{-1}, r^{2}C_{i}r^{-2}, \dots, r^{p-1}C_{i}r^{-(p-1)}$ 构成 $\mathbf{N}$ 的 $p$ 个不同类的那些类。于是 $\mathbf{G}$ 相对于 $\mathbf{N}$ 的类结构确定如下：$\mathbf{N}$ 在 $r$ 下自共轭的类也是 $\mathbf{G}$ 的类；不自共轭的类与 $(p-1)$ 个其他类按上述方式联合，构成 $\mathbf{G}$ 的一个单独的类。

下面说明如何由 $\mathbf{N}$ 的表示确定 $\mathbf{G}$ 的表示。设 $\Delta_{0}$ 是 $\mathbf{N}$ 的维数为 $d$ 的表示，$\phi_{0i}$（$i = 1$ 到 $d$）是 $\Delta_{0}$ 的基，使得对一切 $n \in \mathbf{N}$
\begin{equation*}
n\phi_{0i} = \sum^{d}_{j=1}\phi_{0j}\Delta_{0}(n)_{ji}
\tag{4.5.18}
\end{equation*}
（$i = 1, 2, \dots, d$）。与 4.1 节同样地定义函数
\begin{equation*}
\phi_{mi} = r^{m}\phi_{0i}
\tag{4.5.19}
\end{equation*}
（$m = 0, 1, \dots, (p-1)$，$i = 1, 2, \dots, d$），则类比式 (4.2.16) 可知，对固定的 $m$
\begin{equation*}
n\phi_{mi} = \sum^{d}_{j=1}\phi_{mj}\Delta_{0}(r^{-m}nr^{m})_{ji}
\tag{4.5.20}
\end{equation*}
（$i = 1, 2, \dots, d$）。由 4.2 节的理论可见，这样我们定义了 $\mathbf{N}$ 的 $p$ 个表示 $\Delta_{m}$：
\begin{equation*}
\Delta_{m}(n) = \Delta_{0}(r^{-m}nr^{m})
\tag{4.5.21}
\end{equation*}
（$m = 0, 1, \dots, (p-1)$）。现在 $\mathbf{G}/\mathbf{N}$ 同构于 $p$ 阶循环群 $C_{p}$。由于 $p$ 是素数，只能出现两种情形。$\Delta_{0}$ 的小陪群的阶是 $|\mathbf{G}/\mathbf{N}|$ 的因子，必为 1 或 $p$；$\Delta_{0}$ 的星相应地分别含 $p$ 个表示或 1 个表示。于是得到与类结构类似的结果：要么 $\Delta_{0}$ 在 $r$ 下自共轭、诸 $\Delta_{m}$ 彼此等价、小群是整个群 $\mathbf{G}$；要么诸 $\Delta_{m}$ 彼此不等价、小群恰为不变子群 $\mathbf{N}$。

\noindent\textbf{情形 (i)}\ 设 $\Delta_{m}$（$m = 0, 1, \dots, (p-1)$）彼此不等价。则由 4.2 节的理论立即可知 $\Delta_{0} \uparrow \mathbf{G}$ 不可约。$\Delta_{0} \uparrow \mathbf{G}$ 的特征标由式 (4.2.22) 立刻得到：
\begin{equation*}
\chi_{\Delta_{0} \uparrow \mathbf{G}}(n) = \sum^{p-1}_{m=0}\chi_{\Delta_{0}}(r^{-m}nr^{m}) = \sum^{p-1}_{m=0}\chi_{\Delta_{m}}(n)
\tag{4.5.22}
\end{equation*}
以及
\begin{equation*}
\chi_{\Delta_{0} \uparrow \mathbf{G}}(r^{i}n) = 0
\tag{4.5.23}
\end{equation*}
（$i = 1, 2, \dots, (p-1)$）。

\noindent\textbf{情形 (ii)}\ 设 $\Delta_{m}$（$m = 0, 1, \dots, (p-1)$）彼此等价。则 $\Delta_{0} \uparrow \mathbf{G}$ 可约。由于诸 $\Delta_{m}$ 有相同的特征标，$\Delta_{0} \uparrow \mathbf{G}$ 的特征标由
\begin{equation*}
\chi_{\Delta_{0} \uparrow \mathbf{G}}(n) = p\chi_{\Delta_{0}}(n)
\tag{4.5.24}
\end{equation*}
以及
\begin{equation*}
\chi_{\Delta_{0} \uparrow \mathbf{G}}(r^{i}n) = 0
\tag{4.5.25}
\end{equation*}
（$i = 1, 2, \dots, (p-1)$）给出。于是 $\Delta_{0} \uparrow \mathbf{G}$ 与自身的交织数为
\begin{equation*}
\frac{1}{|\mathbf{G}|}\sum_{g \in \mathbf{G}}\left|\chi_{\Delta_{0} \uparrow \mathbf{G}}(g)\right|^{2} = \frac{1}{|\mathbf{G}|}p^{2}\sum_{n \in \mathbf{N}}\left|\chi_{\Delta_{0}}(n)\right|^{2} = p^{2}|\mathbf{N}|/|\mathbf{G}| = p.
\tag{4.5.26}
\end{equation*}
因此可以想见，此情形下 $\Delta_{0} \uparrow \mathbf{G}$ 可约化为 $p$ 个彼此不等价的 $\mathbf{G}$ 的表示，每个维数为 $d$。下面我们证明确实如此，并特别说明如何实现约化以及如何得到各分表示的特征标。由于 $\mathbf{G}$ 的每个元素都可写成 $r^{i}n$（$i$ 在 0 到 $(p-1)$ 之间，$n \in \mathbf{N}$），为了约化 $\Delta_{0} \uparrow \mathbf{G}$，只需找一个基变换，把 $n$ 与 $r$ 的代表矩阵同时约化到相同的块对角形式。相对于基 $\langle \phi_{0i}, \phi_{1i}, \dots, \phi_{(p-1)i} |$，$n$ 的代表显然是
\begin{equation*}
(\Delta_{0} \uparrow \mathbf{G})(n) = \begin{pmatrix} \Delta_{0}(n) & 0 & 0 & \dots & 0\\ 0 & \Delta_{0}(r^{-1}nr) & 0 & \dots & 0\\ 0 & 0 & \Delta_{0}(r^{-2}nr^{2}) & \dots & 0\\ \vdots & \vdots & \vdots & & \vdots\\ 0 & 0 & 0 & \dots & \Delta_{0}(r^{-p+1}nr^{p-1}) \end{pmatrix}
\tag{4.5.27}
\end{equation*}
又由于 $\phi_{mi} = r^{m}\phi_{0i}$，
\begin{equation*}
\begin{aligned}
r\langle \phi_{0i}, \phi_{1i}, \dots, \phi_{(p-1)i} | &= \langle r\phi_{0i}, r^{2}\phi_{0i}, \dots, r^{p}\phi_{0i} |\\
&= \langle \phi_{1i}, \phi_{2i}, \dots, \phi_{(p-1)i}, \sum_{j}\phi_{0j}\Delta_{0}(r^{p})_{ji} |
\end{aligned}
\tag{4.5.28}
\end{equation*}
其中在最后的分量中用了 $r^{p} \in \mathbf{N}$ 的事实。因此
\begin{equation*}
(\Delta_{0} \uparrow \mathbf{G})(r) = \begin{pmatrix} 0 & 0 & 0 & \dots & 0 & \Delta_{0}(r^{p})\\ \mathbf{1} & 0 & 0 & \dots & 0 & 0\\ 0 & \mathbf{1} & 0 & \dots & 0 & 0\\ \vdots & \vdots & \vdots & & \vdots & \vdots\\ 0 & 0 & 0 & \dots & \mathbf{1} & 0 \end{pmatrix}
\tag{4.5.29}
\end{equation*}
已知 $\Delta_{1}$ 与 $\Delta_{0}$ 等价，故存在幺正矩阵 $\mathbf{P}$ 使得对一切 $n \in \mathbf{N}$
\begin{equation*}
\Delta_{1}(n) = \mathbf{P}\Delta_{0}(n)\mathbf{P}^{-1}.
\tag{4.5.30}
\end{equation*}
于是 $\Delta_{2}(n) = \Delta_{0}(r^{-2}nr^{2}) = \Delta_{0}(r^{-1}n'r) = \Delta_{1}(n') = \mathbf{P}\Delta_{0}(n')\mathbf{P}^{-1} = \mathbf{P}\Delta_{0}(r^{-1}nr)\mathbf{P}^{-1} = \mathbf{P}^{2}\Delta_{0}(n)\mathbf{P}^{-2}$，更一般地用归纳法得
\begin{equation*}
\Delta_{m}(n) = \mathbf{P}^{m}\Delta_{0}(n)\mathbf{P}^{-m}
\tag{4.5.31}
\end{equation*}
（$m = 0, 1, 2, \dots, (p-1)$）。因此定义新基
\begin{equation*}
\psi_{mi} = \sum^{d}_{j=1}\phi_{mj}(\mathbf{P}^{m})_{ji}
\tag{4.5.32}
\end{equation*}
（$m = 0, 1, \dots, (p-1)$，$i = 1, 2, \dots, d$），则简短的计算表明，相对于这个新基
\begin{equation*}
(\Delta_{0} \uparrow \mathbf{G})(n) = \begin{pmatrix} \Delta_{0}(n) & 0 & 0 & \dots & 0\\ 0 & \Delta_{0}(n) & 0 & \dots & 0\\ 0 & 0 & \Delta_{0}(n) & \dots & 0\\ \vdots & \vdots & \vdots & & \vdots\\ 0 & 0 & 0 & \dots & \Delta_{0}(n) \end{pmatrix}
\tag{4.5.33}
\end{equation*}
以及
\begin{equation*}
(\Delta_{0} \uparrow \mathbf{G})(r) = \begin{pmatrix} 0 & 0 & 0 & \dots & 0 & \Delta_{0}(r^{p})\mathbf{P}^{p-1}\\ \mathbf{P}^{-1} & 0 & 0 & \dots & 0 & 0\\ 0 & \mathbf{P}^{-1} & 0 & \dots & 0 & 0\\ \vdots & \vdots & \vdots & & \vdots & \vdots\\ 0 & 0 & 0 & \dots & \mathbf{P}^{-1} & 0 \end{pmatrix}
\tag{4.5.34}
\end{equation*}
下一步证明 $\Delta_{0}(r^{p})\mathbf{P}^{p}$ 是单位矩阵的标量倍。设 $n$ 是 $\mathbf{N}$ 的任一元素，由于 $r^{p} \in \mathbf{N}$，
\begin{equation*}
\begin{aligned}
\Delta_{0}(n)\Delta_{0}(r^{p})\mathbf{P}^{p} &= \Delta_{0}(nr^{p})\mathbf{P}^{p} = \Delta_{0}(r^{p})\Delta_{0}(r^{-p}nr^{p})\mathbf{P}^{p}\\
&= \Delta_{0}(r^{p})\Delta_{1}(r^{-p+1}nr^{p-1})\mathbf{P}^{p} = \Delta_{0}(r^{p})\mathbf{P}\Delta_{0}(r^{-p+1}nr^{p-1})\mathbf{P}^{p-1}\\
&= \Delta_{0}(r^{p})\mathbf{P}^{p}\Delta_{0}(n),
\end{aligned}
\end{equation*}
（重复这一论证）。故 $\Delta_{0}(r^{p})\mathbf{P}^{p}$ 与 $\mathbf{N}$ 的表示 $\Delta_{0}$ 可交换，于是由 Schur 引理
\begin{equation*}
\Delta_{0}(r^{p})\mathbf{P}^{p} = \lambda\mathbf{1}
\tag{4.5.35}
\end{equation*}
其中 $\lambda$ 是一个相位因子。由式 (4.5.30) 可见矩阵 $\mathbf{P}$ 可以乘一个相位因子；特别地可以选取 $\mathbf{P}$ 使 $\lambda = 1$。假定已作了这一选择，并记 $\mathbf{P}^{-1} = \mathbf{R}$，则下列方程成立：
\begin{equation*}
\Delta_{1}(n) = \mathbf{R}^{-1}\Delta_{0}(n)\mathbf{R}
\tag{4.5.36}
\end{equation*}
\begin{equation*}
\Delta_{0}(r^{p}) = \mathbf{R}^{p}
\tag{4.5.37}
\end{equation*}
\begin{equation*}
(\Delta_{0} \uparrow \mathbf{G})(r) = \begin{pmatrix} 0 & 0 & 0 & \dots & 0 & \mathbf{R}\\ \mathbf{R} & 0 & 0 & \dots & 0 & 0\\ 0 & \mathbf{R} & 0 & \dots & 0 & 0\\ \vdots & \vdots & \vdots & & \vdots & \vdots\\ 0 & 0 & 0 & \dots & \mathbf{R} & 0 \end{pmatrix},
\tag{4.5.38}
\end{equation*}
且式 (4.5.33) 对 $(\Delta_{0} \uparrow \mathbf{G})(n)$ 依然成立。张成这个表示的基可以记作 $\langle \chi_{0}, \chi_{1}, \dots, \chi_{p-1} |$，其中 $\chi_{m}$ 代表行矢量 $(\chi_{m1}, \chi_{m2}, \dots, \chi_{md})$。现在寻找进一步的基变换，试图在保持 $(\Delta_{0} \uparrow \mathbf{G})(n)$ 形式不变的同时把 $(\Delta_{0} \uparrow \mathbf{G})(r)$ 块对角化。为看到这确实可能，尝试
\begin{equation*}
\eta = \sum^{p-1}_{m=0}a_{m}\chi_{m}.
\tag{4.5.39}
\end{equation*}
显然，对系数 $a_{m}$ 的任何选择
\begin{equation*}
n\boldsymbol{\eta} = \sum^{p-1}_{m=0}a_{m}n\chi_{m} = \sum^{p-1}_{m=0}a_{m}\chi_{m}\Delta_{0}(n) = \boldsymbol{\eta}\Delta_{0}(n).
\tag{4.5.40}
\end{equation*}
又由式 (4.5.38)，
\begin{equation*}
\begin{aligned}
r\boldsymbol{\eta} &= r(a_{0}\chi_{0} + a_{1}\chi_{1} + \dots + a_{p-1}\chi_{p-1})\\
&= (a_{0}\chi_{1} + a_{1}\chi_{2} + \dots + a_{p-1}\chi_{0})\mathbf{R}
\end{aligned}
\tag{4.5.41}
\end{equation*}
现在的问题是：能否选取 $a_{m}$ 与标量 $\varepsilon$ 使得
\begin{equation*}
r\boldsymbol{\eta} = \boldsymbol{\eta}\varepsilon\mathbf{R}
\tag{4.5.42}
\end{equation*}
比较式 (4.5.41) 与 (4.5.42) 的右端，需要同时满足
\begin{equation*}
\left.
\begin{array}{rl}
\varepsilon a_{0} &= a_{p-1}\\
\varepsilon a_{1} &= a_{0}\\
&\ \vdots\\
\varepsilon a_{p-1} &= a_{p-2}
\end{array}
\right\}
\tag{4.5.43}
\end{equation*}
它容许 $p$ 个可能的解 $\varepsilon = \varepsilon_{q} = \exp(2\pi qi/p)$（$q = 0, 1, \dots, (p-1)$），即 $p$ 次单位根。这意味着张成 $\Delta_{0} \uparrow \mathbf{G}$ 的矢量空间可以约化为 $p$ 个不变子空间，每个维数为 $d$，对应于 $p$ 次单位根的每一个。同样，在由 $\varepsilon = \varepsilon_{q}$ 确定的矢量空间中，有一个矢量 $\boldsymbol{\eta}_{q}$，其系数 $a_{m}$ 由式 (4.5.43) 对应于 $\varepsilon = \varepsilon_{q}$ 的解给出，满足
\begin{equation*}
\left.
\begin{array}{l}
n\boldsymbol{\eta}_{q} = \boldsymbol{\eta}_{q}\Delta_{0}(n)\\
r\boldsymbol{\eta}_{q} = \boldsymbol{\eta}_{q}\varepsilon_{q}\mathbf{R}
\end{array}
\right\}.
\tag{4.5.44}
\end{equation*}
因此我们得到 $\mathbf{G}$ 的 $p$ 个表示，记作 $\Delta_{0q}$，使得
\begin{equation*}
\left.
\begin{array}{l}
\boldsymbol{\Delta}_{0q}(n) = \boldsymbol{\Delta}_{0}(n)\\
\boldsymbol{\Delta}_{0q}(r) = \varepsilon_{q}\mathbf{R}
\end{array}
\right\}
\tag{4.5.45}
\end{equation*}
（$q = 0, 1, \dots, (p-1)$）。最后，由于 $\Delta_{0}$ 在 $\mathbf{N}$ 下不可约，不可能再有进一步的约化，故 $\Delta_{0q}$ 在 $\mathbf{G}$ 下不可约。又，如上面由式 (4.5.26) 所见，$\Delta_{0} \uparrow \mathbf{G}$ 与自身的交织数为 $p$。由于 $\Delta_{0q}$（$q = 0$ 到 $(p-1)$）彼此不等价，这已用尽全部 $p$ 个交织，不允许再有形如 $\Delta_{0q_{1}} \equiv \Delta_{0q_{2}}$（$q_{1} \neq q_{2}$）的等价；也就是说，不可能有任何两个表示 $\Delta_{0q_{1}}$ 与 $\Delta_{0q_{2}}$（$q_{1} \neq q_{2}$）等价。

作为上述理论的应用示例，我们由 $23\ (T)$ 的表示导出点群 $432\ (O)$ 的表示。我们使用表 2.2 与 2.3 中给出的群 $23\ (T)$ 的特征标表与矩阵代表。这里 $\mathbf{G}$ 是群 $432\ (O)$，$\mathbf{N}$ 是群 $23\ (T)$，故 $p = 2$。元素 $r$ 可取为 $C_{2a}$，于是 $r^{p} = C^{2}_{2a} = E$（恒等元）。

先取 $\Delta_{0} = A$（$23\ (T)$ 的）。显然它自共轭，因为对一切 $n \in 23\ (T)$ 有 $A(n) = 1$。由于 $A$ 是一维的且要求 $\mathbf{R}^{2} = 1$，可取 $\mathbf{R} = +1$，于是 $A \uparrow 432\ (O)$ 分解为 $432\ (O)$ 的两个一维表示，即表 2.2 中标记为 $A_{1}$ 与 $A_{2}$ 的。此外，由于单位根是 $+1$ 与 $-1$，有 $A_{1}(C_{2a}) = 1$ 与 $A_{2}(C_{2a}) = -1$。由于对一切 $n \in 23\ (T)$ 有 $A_{1}(n) = A_{2}(n) = A(n)$，这些方程完全确定了 $A_{1}$ 与 $A_{2}$ 的特征标。

再取 $\Delta_{0} = {}^{1}E$（$23\ (T)$ 的）。由于 $C_{2a}C^{+}_{31}C_{2a} = C^{-}_{32}$（见表 1.5），有 $\Delta_{1}(C^{+}_{31}) = \Delta_{0}(C_{2a}C^{+}_{31}C_{2a}) = \Delta_{0}(C^{-}_{32}) = {}^{1}\mathbf{E}(C^{-}_{32}) = \omega^{*} = {}^{2}\mathbf{E}(C^{+}_{31})$。故 $\Delta_{1} = {}^{2}E$。由于 ${}^{1}E$ 与 ${}^{2}E$ 不等价，可知 ${}^{1}E \uparrow 432\ (O)$ 不可约。于是 ${}^{1}E$ 与 ${}^{2}E$ 合起来构成表 2.2 中标记为 $E$ 的 $432\ (O)$ 的表示。表 2.3 中给出的 $E$ 的矩阵与直接由诱导 ${}^{1}E \uparrow 432\ (O)$ 得到的矩阵差一个等价变换。容易证明诱导矩阵由下列三个矩阵生成：
\begin{equation*}
\left.
\begin{array}{l}
\left({}^{1}E \uparrow 432\right)(C^{+}_{31}) = \begin{pmatrix} \omega & 0\\ 0 & \omega^{*} \end{pmatrix}\\[2ex]
\left({}^{1}E \uparrow 432\right)(C_{2x}) = \begin{pmatrix} 1 & 0\\ 0 & 1 \end{pmatrix}\\[2ex]
\left({}^{1}E \uparrow 432\right)(C_{2a}) = \begin{pmatrix} 0 & 1\\ 1 & 0 \end{pmatrix}
\end{array}
\right\}
\tag{4.5.46}
\end{equation*}

最后取 $\Delta_{0} = T$（$23\ (T)$ 的）。$\Delta_{0}$ 的典型矩阵为
\begin{equation*}
\begin{array}{l}
\Delta_{0}(C_{2x}) = \begin{pmatrix} 1 & 0 & 0\\ 0 & -1 & 0\\ 0 & 0 & -1 \end{pmatrix}\\[2ex]
\Delta_{0}(C^{+}_{31}) = \begin{pmatrix} 0 & 0 & 1\\ 1 & 0 & 0\\ 0 & 1 & 0 \end{pmatrix}.
\end{array}
\tag{4.5.47}
\end{equation*}
由于 $C_{2a}C_{2x}C_{2a} = C_{2y}$ 且 $C_{2a}C^{+}_{31}C_{2a} = C^{-}_{32}$，
\begin{equation*}
\begin{array}{l}
\Delta_{1}(C_{2x}) = \Delta_{0}(C_{2y}) = \begin{pmatrix} -1 & 0 & 0\\ 0 & 1 & 0\\ 0 & 0 & -1 \end{pmatrix}\\[2ex]
\Delta_{1}(C^{+}_{31}) = \Delta_{0}(C^{-}_{32}) = \begin{pmatrix} 0 & 1 & 0\\ 0 & 0 & -1\\ -1 & 0 & 0 \end{pmatrix}.
\end{array}
\tag{4.5.48}
\end{equation*}
由此可见 $\Delta_{1}$ 的特征标等于 $\Delta_{0}$ 的特征标，故 $T$ 自共轭。为把 $T \uparrow 432\ (O)$ 约化，需要找一个矩阵 $\mathbf{R}$（见式 (4.5.36) 与 (4.5.37)）使得
\begin{equation*}
\left.
\begin{array}{l}
\mathbf{R}\Delta_{1}(C_{2x}) = \Delta_{0}(C_{2x})\mathbf{R}\\
\mathbf{R}\Delta_{1}(C^{+}_{31}) = \Delta_{0}(C^{+}_{31})\mathbf{R}\\
\mathbf{R}^{2} = \Delta_{0}(C^{2}_{2a}) = \mathbf{1}
\end{array}
\right\}
\tag{4.5.49}
\end{equation*}
利用式 (4.5.47)--(4.5.49) 可以证明满足这些条件的矩阵 $\mathbf{R}$ 是
\begin{equation*}
\mathbf{R} = \begin{pmatrix} 0 & 1 & 0\\ 1 & 0 & 0\\ 0 & 0 & -1 \end{pmatrix}.
\tag{4.5.50}
\end{equation*}
于是 $T \uparrow 432\ (O)$ 分解为 $432\ (O)$ 的两个三维表示 $\Delta_{00}$ 与 $\Delta_{01}$，其矩阵可由式 (4.5.45) 求出。$C_{2a}$ 的矩阵代表则为 $\Delta_{00}(C_{2a}) = +\mathbf{R}$ 与 $\Delta_{01}(C_{2a}) = -\mathbf{R}$，这与表 2.2 中 $T_{1}$、$T_{2}$ 的特征标一致。
